% ============================================================
%  GUÍA 1: ARITMÉTICA CON NÚMEROS NATURALES Y DECIMALES
%  Curso Introductorio de Matemáticas — 1er Año
%  Autor: Ing. Gabriel Astudillo
%  Sesiones cubiertas: Semana 1 (clases 1 a 4)
% ============================================================

\documentclass[12pt, letterpaper]{article}

% ── Paquetes ─────────────────────────────────────────────────

\usepackage{mathwizards-guia}
\mwguia{Guía 1 — Aritmética con Naturales y Decimales}{Ing. Gabriel Astudillo $\cdot$ Curso 1er Año}

\begin{document}
% ============================================================

% ── Portada / Título ─────────────────────────────────────────
\mwportada{Guía de Estudio 1}{Aritmética con Números Naturales y Decimales}{Sumar, Restar, Multiplicar y Dividir como Campeones}

\vspace{4pt}
\begin{cuadroinfo}
  \small
  \textbf{Presentación:} ¡Hola! En esta primera guía vamos a convertirnos en \emph{campeones de la aritmética}. ¿Sabías que cada vez que compras algo en la tiendita, mides tu estatura o repartes dulces con tus amigos, estás usando números con coma (los \textbf{decimales})? Aquí aprenderás a sumarlos, restarlos, multiplicarlos y dividirlos sin miedo y con mucho orden. Cada sección tiene una explicación sencilla, ejemplos resueltos paso a paso y ejercicios que van de lo más fácil a verdaderos retos. ¡Tú puedes!
  \\[4pt]
  \textbf{Nivel de Dificultad:}\quad
  \facil{} Básico / Directo \quad
  \medio{} Intermedio \quad
  \dificil{} Desafiante
\end{cuadroinfo}

\vspace{8pt}

% ============================================================
\section{La Suma}
% ============================================================

\subsection{\textquestiondown Qué son los números naturales y los decimales?}

Los \textbf{números naturales} son los que usamos para \emph{contar} cosas: $0, 1, 2, 3, 4, \dots$ Sirven para responder «¿cuántas canicas tengo?» o «¿cuántos lápices hay en la cartuchera?».

Los \textbf{números decimales} sirven para \emph{medir} cosas que no son cantidades exactas: precios, estaturas, pesos y tiempos. La \textbf{coma} separa la parte entera de la parte decimal:

\begin{cuadroteoria}
  \textbf{Definición:} Un \textbf{número decimal} tiene la forma \emph{parte entera}, coma, \emph{parte decimal}. Por ejemplo, en $3{,}75$ la parte entera es $3$ y la parte decimal es $75$ (que significa $75$ centésimas, o sea $75/100$).
\end{cuadroteoria}

\noindent Algunos ejemplos de la vida real:
\begin{itemize}[leftmargin=*]
  \item Un helado cuesta $2{,}50$ Bs. (la parte decimal son los céntimos)
  \item Ana mide $1{,}42$ metros.
  \item Un paquete de galletas pesa $0{,}25$ kg.
\end{itemize}
\newpage
\subsection{Cómo sumar números decimales}

Sumar es \emph{juntar}. Para sumar decimales solo debes seguir tres pasos:

\begin{cuadropasos}
  \textbf{Pasos para sumar decimales:}
  \begin{enumerate}
    \item \textbf{Alinea las comas} una debajo de la otra (así se alinean enteros con enteros y decimales con decimales).
    \item \textbf{Completa con ceros} a la derecha si un número tiene menos cifras decimales.
    \item \textbf{Suma de derecha a izquierda}, igual que con los naturales, y \textbf{coloca la coma} en el mismo lugar.
  \end{enumerate}
\end{cuadropasos}

\noindent\textbf{Ejemplo resuelto 1:} Calculemos $35{,}7 + 12{,}48$.

Primero completamos con un cero: $35{,}70$. Luego sumamos en columna:

\[
  \begin{array}{rr}
      & 35{,}70 \\
    + & 12{,}48 \\ \hline
      & 48{,}18
  \end{array}
\]

\noindent Respuesta: $35{,}7 + 12{,}48 = 48{,}18$.

\noindent\textbf{Ejemplo resuelto 2:} En la tiendita, Luis compra un cuaderno de $8{,}75$ Bs. y un marcador de $4{,}50$ Bs. \textquestiondown Cuánto gasta?

\[
  8{,}75 + 4{,}50 = 13{,}25
\]

\noindent Luis gasta $13{,}25$ Bs.

\begin{cuadroadvertencia}
  \textbf{Propiedades de la suma} (te sirven para sumar más rápido y para comprobar tus resultados):
  \begin{itemize}
    \item \textbf{Conmutativa:} el orden no cambia el resultado: $a + b = b + a$. Ejemplo: $3 + 5 = 5 + 3 = 8$.
    \item \textbf{Asociativa:} puedes agrupar como quieras: $(a + b) + c = a + (b + c)$. Ejemplo: $(2 + 3) + 4 = 2 + (3 + 4) = 9$.
    \item \textbf{Elemento neutro:} sumar cero no cambia nada: $a + 0 = a$. Ejemplo: $12{,}5 + 0 = 12{,}5$.
  \end{itemize}
\end{cuadroadvertencia}

\subsection{Ejercicios: La Suma}

\begin{enumerate}[leftmargin=*]
  \item \facil{} $348 + 217$

  \item \facil{} $56 + 89$

  \item \facil{} $1250 + 3375$

  \item \facil{} $3{,}6 + 2{,}5$ % lleva 1 en las décimas

  \item \facil{} $12{,}3 + 45{,}6$

  \item \facil{} $0{,}7 + 0{,}2$

  \item \facil{} $4{,}6 + 3{,}8$ % lleva 1 en las décimas

  \item \facil{} $0{,}9 + 0{,}7$ % lleva 1 a la unidad

  \item \medio{} $458{,}75 + 239{,}64$

  \item \medio{} $2\,458{,}65 + 719{,}834$

  \item \medio{} $45\,892 + 12\,608$

  \item \medio{} $14{,}085 + 230{,}4 + 1\,589{,}76$

  \item \medio{} Ana compró un televisor de $148{,}50$ Bs. y una lavadora de $325{,}75$ Bs. \textquestiondown Cuánto gastó en total?

  \item \medio{} $1\,425{,}8 + 368{,}47 + 79{,}055$

  \item \medio{} Tres amigos juntan su dinero para un proyecto: uno pone $125{,}50$ Bs., otro $88{,}75$ Bs. y el tercero $215{,}25$ Bs. \textquestiondown Cuánto reunieron?

  \item \medio{} Encuentra el número que falta: $456{,}8 + \_\_\_ = 1\,000$

  \item \dificil{} En una tienda de electrodomésticos, Sofía compra una licuadora de $48{,}25$ Bs., una cafetera de $136{,}60$ Bs. y un microondas de $257{,}15$ Bs. Si paga con un billete de $500$ Bs., \textquestiondown cuánto recibe de vuelto? (Pista: primero suma todo, luego resta.)

  \item \dificil{} $145{,}125 + 234{,}875 + 1\,620{,}5$

  \item \dificil{} Un atleta entrena en tres etapas: en la primera recorre $14{,}25$ km, en la segunda $28{,}5$ km y en la tercera $3\,750$ m. \textquestiondown Cuántos kilómetros recorre en total? (Recuerda: $3\,750$ m $= 3{,}75$ km.)

  \item \dificil{} \textbf{El cuadrado mágico.} Coloca los números $1{,}1;\ 1{,}2;\ 1{,}3;\ \dots;\ 1{,}9$ (todos una vez) en el cuadrado de $3 \times 3$ de manera que la suma de cada fila, de cada columna y de cada diagonal sea siempre $4{,}5$. \textquestiondown Puedes lograrlo?

  \item \dificil{} Un edificio mide $148{,}5$ m de alto. Le agregan una torre de transmisión de $27{,}85$ m y encima una antena de $4{,}65$ m. \textquestiondown Cuánto mide ahora la estructura completa?

  \item \dificil{} \textbf{:} Calcula la suma agrupando los números convenientemente (esa es la pista): $99{,}9 + 0{,}1 + 87{,}5 + 12{,}5 + 775{,}5 + 224{,}5$
\end{enumerate}

\newpage

% ============================================================
\section{La Resta}
% ============================================================

\subsection{Cómo restar números decimales}

Restar es \emph{quitar}. Para restar decimales usamos los mismos pasos que para sumar: \textbf{alinear las comas}, \textbf{completar con ceros} y \textbf{restar de derecha a izquierda}. Si una cifra no alcanza, \emph{pedimos prestado} a la izquierda, igual que con los naturales.

\begin{cuadropasos}
  \textbf{Paso extra (¡importante!):} después de restar, \textbf{comprueba} tu respuesta sumando: si $a - b = c$, entonces debe cumplirse que $c + b = a$. Así sabrás si tu resta está bien.
\end{cuadropasos}

\noindent\textbf{Ejemplo resuelto 1:} Calculemos $100 - 37{,}25$.

Completamos con ceros: $100{,}00 - 37{,}25$.

\[
  \begin{array}{rr}
      & 100{,}00 \\
    - & 37{,}25  \\ \hline
      & 62{,}75
  \end{array}
\]

\noindent Comprobación: $62{,}75 + 37{,}25 = 100{,}00$. \textquestiondown Correcto!

\noindent\textbf{Ejemplo resuelto 2:} Tenías $20$ Bs. y gastaste $12{,}75$ Bs. en el kiosco. \textquestiondown Cuánto te queda?

\[
  20{,}00 - 12{,}75 = 7{,}25
\]

\noindent Te quedan $7{,}25$ Bs.

\begin{cuadroadvertencia}
  \textbf{Ojo con esto:} en la resta el orden \textbf{sí} importa. $8 - 3 = 5$, pero $3 - 8$ no da $5$ (en naturales no se puede quitar más de lo que hay... ¡pero en la guía 2 conocerás los números negativos que lo resuelven!).
\end{cuadroadvertencia}

\subsection{Ejercicios: La Resta}

\begin{enumerate}[leftmargin=*]
  \item \facil{} $789 - 456$

  \item \facil{} $1000 - 347$

  \item \facil{} $93 - 57$

  \item \facil{} $72 - 38$ % pide prestado (simple)

  \item \facil{} $200 - 46$ % pide prestado (recursivo: pasa por el cero)

  \item \facil{} $8{,}9 - 3{,}4$

  \item \facil{} $15{,}6 - 12{,}3$

  \item \facil{} $1 - 0{,}5$

  \item \medio{} $412{,}4 - 185{,}75$

  \item \medio{} $2\,500 - 1\,347{,}85$

  \item \medio{} $650{,}03 - 278{,}48$

  \item \medio{} Tenías $500$ Bs. y gastaste $348{,}65$ Bs. \textquestiondown Cuánto te queda?

  \item \medio{} $100 - 0{,}999$

  \item \medio{} Encuentra el número que falta: $\_\_\_ - 145{,}5 = 354{,}5$

  \item \medio{} Un rascacielos mide $175{,}8$ m y un edificio vecino mide $142{,}35$ m. \textquestiondown Cuál es la diferencia de altura entre ambos?

  \item \medio{} $4\,500{,}005 - 1\,234{,}56$

  \item \dificil{} María compra tres artefactos que cuestan $125{,}5$ Bs., $88{,}75$ Bs. y $166{,}25$ Bs. Si paga con un billete de $500$ Bs., \textquestiondown cuánto recibe de vuelto?

  \item \dificil{} Durante el día la temperatura subió $18{,}5\,^{\circ}$C y luego por la noche bajó $13{,}75\,^{\circ}$C. \textquestiondown Cuál fue la variación neta de temperatura?

  \item \dificil{} $30 - 0{,}3 - 0{,}03 - 0{,}003$

  \item \dificil{} $500 - 135{,}5 - 214{,}25 - 80{,}25$

  \item \dificil{} Un depósito contiene $255{,}5$ litros de agua. Se extraen tres recipientes de $35{,}5$ litros cada uno. \textquestiondown Cuánta agua queda en el depósito?

  \item \dificil{} \textbf{:} Calcula con orden: $100 - 99{,}9 + 88{,}8 - 77{,}7 + 66{,}6 - 55{,}5$. (Hazlo paso a paso, de izquierda a derecha.)
\end{enumerate}


% ============================================================
\section{La Multiplicación}
% ============================================================

\subsection{Cómo multiplicar números decimales}

Multiplicar es \emph{sumar varias veces lo mismo}. Por ejemplo, $4 \times 3$ es lo mismo que $3 + 3 + 3 + 3 = 12$. Por eso decimos que la multiplicación es una \textbf{suma repetida}.

Multiplicar decimales es muy fácil si recuerdas esta idea:

\begin{cuadropasos}
  \textbf{Pasos para multiplicar decimales:}
  \begin{enumerate}
    \item \textbf{Olvídate de la coma} por un momento y multiplica los números como si fueran naturales.
    \item \textbf{Cuenta las cifras decimales} de los dos números y súmalas.
    \item En el resultado, \textbf{cuenta esa misma cantidad de cifras de derecha a izquierda} y coloca la coma ahí.
  \end{enumerate}
\end{cuadropasos}

\noindent\textbf{Ejemplo resuelto 1:} Calculemos $2{,}5 \times 1{,}3$.

Multiplicamos sin coma: $25 \times 13 = 325$. Ahora contamos cifras decimales: $2{,}5$ tiene una cifra decimal y $1{,}3$ tiene una cifra decimal. En total: $2$ cifras. Colocamos la coma: $3{,}25$.

\noindent Respuesta: $2{,}5 \times 1{,}3 = 3{,}25$.

\noindent\textbf{Ejemplo resuelto 2:} \textquestiondown Cuánto cuestan $4$ helados de $3{,}50$ Bs. cada uno?

\[
  4 \times 3{,}50 = 14{,}00
\]

\noindent Los cuatro helados cuestan $14$ Bs.

\begin{cuadroadvertencia}
  \textbf{Propiedades de la multiplicación:}
  \begin{itemize}
    \item \textbf{Conmutativa:} $a \times b = b \times a$. Ejemplo: $3 \times 4 = 4 \times 3 = 12$.
    \item \textbf{Asociativa:} $(a \times b) \times c = a \times (b \times c)$.
    \item \textbf{Distributiva:} $a \times (b + c) = a \times b + a \times c$. Ejemplo: $2 \times (3 + 4) = 2 \times 3 + 2 \times 4 = 6 + 8 = 14$.
    \item \textbf{Elemento neutro:} multiplicar por $1$ no cambia nada: $a \times 1 = a$.
    \item \textbf{Elemento absorbente:} multiplicar por $0$ da $0$: $a \times 0 = 0$.
  \end{itemize}
\end{cuadroadvertencia}

\subsection{Ejercicios: La Multiplicación}

\begin{enumerate}[leftmargin=*]
  \item \facil{} $23 \times 4$

  \item \facil{} $305 \times 2$

  \item \facil{} $12 \times 11$

  \item \facil{} $3 \times 0{,}5$

  \item \facil{} $0{,}4 \times 5$

  \item \facil{} $2 \times 2{,}5$

  \item \medio{} $1\,250 \times 345$

  \item \medio{} $4\,350 \times 128$

  \item \medio{} $6\,250 \times 164$

  \item \medio{} $14\,250 \times 35$

  \item \medio{} $2\,250 \times 140$

  \item \medio{} En una tienda hay $240$ paquetes de galletas y cada uno cuesta $13{,}25$ Bs. \textquestiondown Cuánto cuestan todos juntos?

  \item \medio{} $1\,420{,}5 \times 35{,}4$

  \item \medio{} $650{,}8 \times 42{,}5$

  \item \medio{} $4\,500 \times 37{,}5$

  \item \dificil{} $12\,500 \times 348$

  \item \dificil{} Un terreno rectangular mide $450{,}5$ m de largo y $230{,}4$ m de ancho. \textquestiondown Cuál es su área? (Recuerda: área $=$ largo $\times$ ancho.)

  \item \dificil{} $0{,}125 \times 48\,000$ y también $0{,}0625 \times 32\,000$

  \item \dificil{} \textquestiondown Cuánto cuestan $1\,200$ botellas de jugo si cada una cuesta $8{,}50$ Bs.?
  \item \dificil{} \textbf{:}
        \begin{enumerate}[label=\alph*)]
          \item Calcula $250 \times 125 \times 40$ agrupando primero $250 \times 40$.
          \item David ahorra $125{,}50$ Bs. cada semana. \textquestiondown Cuánto habrá ahorrado en $480$ semanas?
        \end{enumerate}
\end{enumerate}

% ============================================================
\section{La Unidad Seguida de Ceros}
% ============================================================

\subsection{Cómo multiplicar por 10, 100 y 1000}

Los números $10$, $100$ y $1000$ se llaman \textbf{unidades seguidas de ceros}. Multiplicar por ellos es facilísimo: \textbf{no hay que multiplicar}, ¡solo se mueve la coma!

\begin{cuadrosolucion}
  \textbf{Multiplicar por $10$, $100$ o $1000$:} mueve la coma hacia la \textbf{derecha} $1$, $2$ o $3$ lugares (la misma cantidad de ceros que tenga el número). Si se acaban las cifras, \textbf{agrega ceros}.
  \begin{itemize}
    \item $4{,}25 \times 10 = 42{,}5$ (la coma se movió 1 lugar)
    \item $0{,}7 \times 100 = 70$ (la coma se movió 2 lugares y agregamos un cero)
    \item $35 \times 1000 = 35\,000$ (agregamos 3 ceros)
  \end{itemize}
\end{cuadrosolucion}

\noindent\textbf{Ejemplo resuelto 1:} Una moneda pesa $0{,}0035$ kg. \textquestiondown Cuánto pesan $100$ monedas?

\[
  0{,}0035 \times 100 = 0{,}35
\]

\noindent Las $100$ monedas pesan $0{,}35$ kg.
\newpage
\subsection{Cómo dividir por 10, 100 y 1000}
Dividir por $10$, $100$ o $1000$ es el \emph{camino inverso}: ahora la coma se mueve hacia la \textbf{izquierda}.

\begin{cuadroteoria}
  \textbf{Dividir por $10$, $100$ o $1000$:} mueve la coma hacia la \textbf{izquierda} $1$, $2$ o $3$ lugares. Si se acaban las cifras, \textbf{agrega ceros} delante.
  \begin{itemize}
    \item $42{,}5 \div 10 = 4{,}25$
    \item $700 \div 100 = 7$
    \item $3 \div 1000 = 0{,}003$ (agregamos dos ceros delante)
  \end{itemize}
\end{cuadroteoria}

\noindent\textbf{Ejemplo resuelto 2:} Un paquete de $10$ bolígrafos cuesta $15{,}50$ Bs. \textquestiondown Cuánto cuesta cada bolígrafo?

\[
  15{,}50 \div 10 = 1{,}55
\]

\noindent Cada bolígrafo cuesta $1{,}55$ Bs.

\begin{cuadroadvertencia}
  \textbf{¿Para qué nos sirve?} ¡Es la llave de las divisiones con decimales! Cuando el divisor tiene coma, \textbf{multiplicamos los dos números por $10$, $100$ o $1000$} para quitarle la coma:
  \[
    3 \div 0{,}5 = 30 \div 5 = 6 \qquad \text{(multiplicamos por $10$)}
  \]
  \[
    0{,}6 \div 0{,}12 = 60 \div 12 = 5 \qquad \text{(multiplicamos por $100$)}
  \]
  Dominar la unidad seguida de ceros te hará dividir muchísimo más rápido.
\end{cuadroadvertencia}

\subsection{Ejercicios: La Unidad Seguida de Ceros}

\begin{enumerate}[leftmargin=*]
  \item \facil{} $7 \times 10$

  \item \facil{} $34 \times 10$

  \item \facil{} $5 \times 100$

  \item \facil{} $2{,}5 \times 10$

  \item \facil{} $4{,}75 \times 100$

  \item \facil{} $80 \div 10$

  \item \facil{} $700 \div 100$

  \item \facil{} $3{,}5 \div 10$

  \item \medio{} $0{,}8 \times 10$

  \item \medio{} $0{,}125 \times 1000$

  \item \medio{} $56 \times 1000$

  \item \medio{} $450 \div 10$

  \item \medio{} $12{,}5 \div 10$

  \item \medio{} $2\,500 \div 100$

  \item \medio{} $0{,}7 \div 10$

  \item \medio{} $4\,800 \div 1000$

  \item \medio{} Un paquete de $100$ tornillos pesa $0{,}85$ kg. \textquestiondown Cuánto pesa cada tornillo?

  \item \dificil{} $0{,}04 \times 100$

  \item \dificil{} $1{,}5 \times 1000$

  \item \dificil{} \textbf{:}
        \begin{enumerate}[label=\alph*)]
          \item Calcula $2{,}5 \times 1000 \div 100$ (primero multiplica, luego divide).
          \item \textquestiondown Cuántas veces hay que dividir $32\,000$ entre $10$ para obtener $0{,}32$?
        \end{enumerate}
\end{enumerate}

% ============================================================
\section{La División}
% ============================================================

\subsection{Cómo dividir con decimales}

Dividir es \emph{repartir en partes iguales}. Si repartimos $12$ caramelos entre $4$ amigos, cada uno recibe $12 \div 4 = 3$.

Cuando aparece una coma en la división hay tres casos:

\begin{cuadroteoria}
  \textbf{Caso 1: Decimal entre entero.} Divide normalmente. Cuando bajes la coma del dividendo, coloca la coma en el cociente.
  \[
    4{,}8 \div 4 = 1{,}2
  \]

  \textbf{Caso 2: Entero entre decimal.} Multiplica \emph{los dos números} por $10$, $100$ o $1000$ (lo que haga falta, como viste en la sección anterior) para que el divisor quede sin coma. Luego divide como enteros.
  \[
    3 \div 0{,}5 = 30 \div 5 = 6
  \]

  \textbf{Caso 3: Decimal entre decimal.} Igual que el caso 2: mueve la coma del divisor hasta que sea entero y mueve la del dividendo la misma cantidad de lugares.
  \[
    0{,}6 \div 0{,}12 = 60 \div 12 = 5
  \]
\end{cuadroteoria}
\newpage
\begin{cuadrosolucion}
  \textbf{Comprobación de la división:} si la división es exacta, debe cumplirse que
  \[
    \text{divisor} \times \text{cociente} = \text{dividendo}.
  \]
  Si sobra residuo: $\text{divisor} \times \text{cociente} + \text{residuo} = \text{dividendo}$.
\end{cuadrosolucion}

\begin{cuadroadvertencia}
  \textbf{Divisiones inexactas y cifras decimales:} \\
  Una división es \textbf{inexacta} cuando el resto no llega a ser cero, aunque sigamos sacando decimales. \\
  \textbf{Regla para esta guía:} Cuando realices una división inexacta con decimales, la calcularás \textbf{únicamente hasta dos cifras decimales (centésimas)}.
\end{cuadroadvertencia}

% ── El método de la galera ──────────────────────────────────
\subsection{El método de la galera (división paso a paso)}

Para dividir usamos el \textbf{método de la galera}: colocamos el dividendo \emph{dentro} de la galera y el divisor \emph{fuera, a la derecha}. El cociente se escribe \emph{debajo} del divisor.

\begin{cuadropasos}
  \textbf{Pasos del método de la galera:}
  \begin{enumerate}
    \item Escribe el \textbf{dividendo} dentro de la galera y el \textbf{divisor} a la derecha, separados por una línea vertical y una línea horizontal debajo del divisor.
    \item Pregúntate: \textquestiondown cuántas veces cabe el divisor en las primeras cifras del dividendo?
    \item Escribe ese número en el \textbf{cociente} (debajo del divisor).
    \item \textbf{Multiplica} ese dígito por el divisor y \textbf{resta} el producto al grupo de cifras que estás dividiendo.
    \item \textbf{Baja} la siguiente cifra del dividendo al lado del resto.
    \item Repite los pasos 2 a 5 hasta que no queden cifras por bajar.
    \item Si al terminar queda un resto y necesitas decimales: coloca la \textbf{coma} en el cociente, agrega un \textbf{cero} al resto y continúa dividiendo.
  \end{enumerate}
\end{cuadropasos}
\newpage
% ── Caso 1: Decimal entre Entero ────────────────────────────
\subsubsection*{Caso 1: Decimal entre entero}

\noindent\textbf{Ejemplo resuelto 1:} Calculemos $\boldsymbol{48{,}6 \div 3}$ con la galera.

\begin{cuadrosolucion}
  \textbf{Paso 1:} Armamos la galera:

  \[
    \renewcommand{\arraystretch}{1.2}
    \begin{array}{r|l}
      48{,}6 & 3 \\ \hline
             &
    \end{array}
  \]

  \textbf{Paso 2:} \textquestiondown Cuántas veces cabe $3$ en $4$? Cabe $1$ vez ($1 \times 3 = 3$). Escribimos $1$ en el cociente. Restamos: $4 - 3 = 1$.

  \[
    \renewcommand{\arraystretch}{1.2}
    \begin{array}{r|l}
      4\underline{8}{,}6         & 3 \\ \hline
      \mathbf{-3}\phantom{8{,}6} & 1 \\
      1\phantom{8{,}6}           &
    \end{array}
  \]

  \textbf{Paso 3:} Bajamos el $8$. Ahora tenemos $18$. \textquestiondown Cuántas veces cabe $3$ en $18$? Cabe $6$ veces ($6 \times 3 = 18$). Restamos: $18 - 18 = 0$.

  \[
    \renewcommand{\arraystretch}{1.2}
    \begin{array}{r|l}
      48{,}\underline{6}         & 3  \\ \hline
      \mathbf{-3}\phantom{8{,}6} & 16 \\
      1\mathbf{8}\phantom{{,}6}  &    \\
      \mathbf{-18}\phantom{{,}6} &    \\
      0\phantom{{,}6}            &
    \end{array}
  \]

  \textbf{Paso 4:} Bajamos el $6$. Como pasamos la coma en el dividendo, \textbf{colocamos la coma en el cociente}: $16{,}\ldots$ Ahora tenemos $06$. \textquestiondown Cuántas veces cabe $3$ en $6$? Cabe $2$ veces ($2 \times 3 = 6$). Restamos: $6 - 6 = 0$.

  \[
    \renewcommand{\arraystretch}{1.2}
    \begin{array}{r|l}
      48{,}6 & 3                   \\ \hline
             & \boldsymbol{16{,}2} \\
    \end{array}
    \qquad \text{Resto } = 0
  \]

  \textbf{Respuesta:} $48{,}6 \div 3 = \boldsymbol{16{,}2}$.\\
  \textbf{Comprobación:} $16{,}2 \times 3 = 48{,}6$ \checkmark
\end{cuadrosolucion}

\newpage %ajuste

% ── Caso 2: Entero entre Decimal ────────────────────────────
\subsubsection*{Caso 2: Entero entre decimal}

\noindent\textbf{Ejemplo resuelto 2:} Calculemos $\boldsymbol{45 \div 0{,}5}$ con la galera.

\begin{cuadrosolucion}
  \textbf{Preparación:} El divisor $0{,}5$ tiene una cifra decimal. Multiplicamos ambos por $10$ para quitarle la coma:
  \[
    45 \div 0{,}5 \;=\; 450 \div 5
  \]

  Ahora dividimos $450 \div 5$ con la galera:

  \textbf{Paso 1:} \textquestiondown Cuántas veces cabe $5$ en $4$? No cabe ($4 < 5$). Tomamos las dos primeras cifras: $45$. \textquestiondown Cuántas veces cabe $5$ en $45$? Cabe $9$ veces ($9 \times 5 = 45$). Restamos: $45 - 45 = 0$.

  \[
    \renewcommand{\arraystretch}{1.2}
    \begin{array}{r|l}
      45\underline{0}         & 5 \\ \hline
      \mathbf{-45}\phantom{0} & 9 \\
      0\phantom{0}            &
    \end{array}
  \]

  \textbf{Paso 2:} Bajamos el $0$. Tenemos $00$. \textquestiondown Cuántas veces cabe $5$ en $0$? Cabe $0$ veces.

  \[
    \renewcommand{\arraystretch}{1.2}
    \begin{array}{r|l}
      450 & 5               \\ \hline
          & \boldsymbol{90} \\
    \end{array}
    \qquad \text{Resto } = 0
  \]

  \textbf{Respuesta:} $45 \div 0{,}5 = \boldsymbol{90}$.\\
  \textbf{Comprobación:} $90 \times 0{,}5 = 45$ \checkmark
\end{cuadrosolucion}

% ── Caso 3: Decimal entre Decimal ────────────────────────────
\subsubsection*{Caso 3: Decimal entre decimal}

\noindent\textbf{Ejemplo resuelto 3:} Calculemos $\boldsymbol{7{,}56 \div 0{,}12}$ con la galera.

\begin{cuadrosolucion}
  \textbf{Preparación:} El divisor $0{,}12$ tiene dos cifras decimales. Multiplicamos ambos por $100$:
  \[
    7{,}56 \div 0{,}12 \;=\; 756 \div 12
  \]

  Ahora dividimos $756 \div 12$ con la galera:

  \textbf{Paso 1:} \textquestiondown Cuántas veces cabe $12$ en $7$? No cabe. Tomamos $75$. \textquestiondown Cuántas veces cabe $12$ en $75$? Probamos: $6 \times 12 = 72$ (cabe). $7 \times 12 = 84$ (no cabe, es mayor que $75$). Usamos $6$. Restamos: $75 - 72 = 3$.

  \[
    \renewcommand{\arraystretch}{1.2}
    \begin{array}{r|l}
      75\underline{6}         & 12 \\ \hline
      \mathbf{-72}\phantom{6} & 6  \\
      3\phantom{6}            &
    \end{array}
  \]

  \textbf{Paso 2:} Bajamos el $6$. Tenemos $36$. \textquestiondown Cuántas veces cabe $12$ en $36$? Cabe $3$ veces ($3 \times 12 = 36$). Restamos: $36 - 36 = 0$.

  \[
    \renewcommand{\arraystretch}{1.2}
    \begin{array}{r|l}
      756 & 12              \\ \hline
          & \boldsymbol{63} \\
    \end{array}
    \qquad \text{Resto } = 0
  \]

  \textbf{Respuesta:} $7{,}56 \div 0{,}12 = \boldsymbol{63}$.\\
  \textbf{Comprobación:} $63 \times 0{,}12 = 7{,}56$ \checkmark
\end{cuadrosolucion}

% ── División inexacta ────────────────────────────────────────
\subsubsection*{División inexacta (con decimales en el cociente)}

\noindent\textbf{Ejemplo resuelto 4:} Calculemos $\boldsymbol{14 \div 3}$ hasta dos cifras decimales con la galera.

\begin{cuadrosolucion}
  \textbf{Paso 1:} \textquestiondown Cuántas veces cabe $3$ en $1$? No cabe. Tomamos $14$. \textquestiondown Cuántas veces cabe $3$ en $14$? Cabe $4$ veces ($4 \times 3 = 12$). Restamos: $14 - 12 = 2$.

  \[
    \renewcommand{\arraystretch}{1.2}
    \begin{array}{r|l}
      14           & 3 \\ \hline
      \mathbf{-12} & 4 \\
      2            &
    \end{array}
  \]

  \textbf{Paso 2:} No quedan más cifras, pero queda resto. Colocamos la \textbf{coma} en el cociente y agregamos un cero al resto: tenemos $20$. \textquestiondown Cuántas veces cabe $3$ en $20$? Cabe $6$ veces ($6 \times 3 = 18$). Restamos: $20 - 18 = 2$.

  \[
    \renewcommand{\arraystretch}{1.2}
    \begin{array}{r|l}
      14{,}0                     & 3     \\ \hline
      \mathbf{-12}\phantom{{,}0} & 4{,}6 \\
      2\mathbf{0}\phantom{}      &       \\
      \mathbf{-18}\phantom{}     &       \\
      2\phantom{}                &
    \end{array}
  \]

  \textbf{Paso 3:} Agregamos otro cero: tenemos $20$ de nuevo. Cabe $6$ veces otra vez. Restamos: $20 - 18 = 2$.

  Ya alcanzamos las \textbf{dos cifras decimales} (centésimas), así que nos detenemos.

  \[
    \renewcommand{\arraystretch}{1.2}
    \begin{array}{r|l}
      14{,}00 & 3                   \\ \hline
              & \boldsymbol{4{,}66} \\
    \end{array}
    \qquad \text{Resto } = 0{,}02
  \]

  \textbf{Respuesta:} $14 \div 3 = \boldsymbol{4{,}66}$ (con resto $0{,}02$).\\
  \textbf{Comprobación:} $4{,}66 \times 3 + 0{,}02 = 13{,}98 + 0{,}02 = 14$ \checkmark
\end{cuadrosolucion}

\newpage
% ── Problema de la vida real ─────────────────────────────────
\subsubsection*{¡Otro ejemplo!}

\noindent\textbf{Ejemplo resuelto 5:} Repartimos $\boldsymbol{27{,}60}$ Bs. entre $\boldsymbol{3}$ amigos. \textquestiondown Cuánto le toca a cada uno?

\begin{cuadrosolucion}
  Dividimos $27{,}60 \div 3$ con la galera:

  \textbf{Paso 1:} \textquestiondown Cuántas veces cabe $3$ en $2$? No cabe. Tomamos $27$. \textquestiondown Cuántas veces cabe $3$ en $27$? Cabe $9$ veces ($9 \times 3 = 27$). Restamos: $27 - 27 = 0$.

  \[
    \renewcommand{\arraystretch}{1.2}
    \begin{array}{r|l}
      27{,}\underline{60}         & 3 \\ \hline
      \mathbf{-27}\phantom{{,}60} & 9 \\
      0\phantom{{,}60}            &
    \end{array}
  \]

  \textbf{Paso 2:} Bajamos el $6$. Como pasamos la coma, \textbf{colocamos la coma en el cociente}: $9{,}\ldots$ Tenemos $06$. \textquestiondown Cuántas veces cabe $3$ en $6$? Cabe $2$ veces ($2 \times 3 = 6$). Restamos: $6 - 6 = 0$.

  \textbf{Paso 3:} Bajamos el $0$. Tenemos $00$. \textquestiondown Cuántas veces cabe $3$ en $0$? Cabe $0$ veces.

  \[
    \renewcommand{\arraystretch}{1.2}
    \begin{array}{r|l}
      27{,}60 & 3                   \\ \hline
              & \boldsymbol{9{,}20} \\
    \end{array}
    \qquad \text{Resto } = 0
  \]

  \textbf{Respuesta:} Cada amigo recibe $\boldsymbol{9{,}20}$ Bs.\\
  \textbf{Comprobación:} $9{,}20 \times 3 = 27{,}60$ \checkmark
\end{cuadrosolucion}

\subsection{Ejercicios: La División}

\begin{enumerate}[leftmargin=*]
  \item \facil{} $84 \div 4$

  \item \facil{} $97 \div 3$

  \item \facil{} $125 \div 5$

  \item \facil{} $19 \div 3$

  \item \facil{} $7{,}3 \div 3$

  \item \facil{} $2{,}4 \div 2$

  \item \medio{} $156\,250 \div 25$

  \item \medio{} $486\,500 \div 12$

  \item \medio{} $125{,}4 \div 0{,}7$

  \item \medio{} $1\,562\,500 \div 125$

  \item \medio{} $812{,}5 \div 0{,}3$

  \item \medio{} $435\,000 \div 150$

  \item \medio{} $15\,625\,000 \div 1\,250$

  \item \medio{} $7\,680\,000 \div 3\,200$

  \item \medio{} $8\,125\,000 \div 2\,500$

  \item \dificil{} $15\,625 \div 0{,}125$

  \item \dificil{} $450 \div 0{,}7$

  \item \dificil{} Tres hermanos se reparten en partes iguales un premio de $27{,}60$ Bs. \textquestiondown Cuánto recibe cada uno?

  \item \dificil{} Un rollo de tela de $25{,}5$ m se usa para confeccionar $7$ vestidos iguales. \textquestiondown Cuántos metros de tela se emplean por vestido? (Calcula el resultado hasta dos decimales).

  \item \dificil{} \textbf{:}
        \begin{enumerate}[label=\alph*)]
          \item Calcula $(3\,600 \div 45) \div (0{,}25 \times 16)$.
          \item Pienso un número: si lo divido entre $0{,}25$ obtengo $40$. \textquestiondown Cuál es el número? (Trabaja hacia atrás.)
        \end{enumerate}
\end{enumerate}

\newpage

% ============================================================
\ifmwclaves
  \section{Clave de Respuestas}
  % ============================================================

  \subsection*{La Suma (Ejercicios 1--22)}

  \begin{enumerate}[leftmargin=*, label=\textbf{\arabic*.}]
    \item $565$
    \item $145$
    \item $4625$
    \item $6{,}1$
    \item $57{,}9$
    \item $0{,}9$
    \item $8{,}4$
    \item $1{,}6$
    \item $458{,}75 + 239{,}64 = 698{,}39$
    \item $2\,458{,}65 + 719{,}834 = 3\,178{,}484$
    \item $45\,892 + 12\,608 = 58\,500$
    \item $14{,}085 + 230{,}4 + 1\,589{,}76 = 1\,834{,}245$
    \item $148{,}50 + 325{,}75 = 474{,}25$ Bs.
    \item $1\,425{,}8 + 368{,}47 + 79{,}055 = 1\,873{,}325$
    \item $125{,}50 + 88{,}75 + 215{,}25 = 429{,}50$ Bs.
    \item $456{,}8 + 543{,}2 = 1\,000$
    \item Total: $48{,}25 + 136{,}60 + 257{,}15 = 442{,}00$ Bs. Vuelto: $500 - 442 = 58$ Bs.
    \item $145{,}125 + 234{,}875 = 380$; luego $380 + 1\,620{,}5 = 2\,000{,}5$
    \item $14{,}25 + 28{,}5 + 3{,}75 = 46{,}5$ km
    \item Una posible solución (todas las filas, columnas y diagonales suman $4{,}5$):
          \[
            \begin{array}{ccc}
              1{,}6 & 1{,}1 & 1{,}8 \\
              1{,}7 & 1{,}5 & 1{,}3 \\
              1{,}2 & 1{,}9 & 1{,}4
            \end{array}
          \]
    \item $148{,}5 + 27{,}85 + 4{,}65 = 181{,}00$ m
    \item Agrupando: $(99{,}9 + 0{,}1) + (87{,}5 + 12{,}5) + (775{,}5 + 224{,}5) = 100 + 100 + 1000 = 1\,200$
  \end{enumerate}

  \subsection*{La Resta (Ejercicios 23--44)}

  \begin{enumerate}[leftmargin=*, label=\textbf{\arabic*.}]
    \setcounter{enumi}{22}
    \item $333$
    \item $653$
    \item $36$
    \item $34$
    \item $154$
    \item $5{,}5$
    \item $3{,}3$
    \item $0{,}5$
    \item $412{,}4 - 185{,}75 = 226{,}65$
    \item $2\,500 - 1\,347{,}85 = 1\,152{,}15$
    \item $650{,}03 - 278{,}48 = 371{,}55$
    \item $500 - 348{,}65 = 151{,}35$ Bs.
    \item $100 - 0{,}999 = 99{,}001$
    \item $500 - 145{,}5 = 354{,}5$
    \item $175{,}8 - 142{,}35 = 33{,}45$ m
    \item $4\,500{,}005 - 1\,234{,}56 = 3\,265{,}445$
    \item Total: $125{,}5 + 88{,}75 + 166{,}25 = 380{,}50$ Bs. Vuelto: $500 - 380{,}50 = 119{,}50$ Bs.
    \item Subió $18{,}5$ y bajó $13{,}75$: $18{,}5 - 13{,}75 = 4{,}75\,^{\circ}$C.
    \item $30 - 0{,}3 = 29{,}7$; $- 0{,}03 = 29{,}67$; $- 0{,}003 = 29{,}667$
    \item $500 - 135{,}5 - 214{,}25 - 80{,}25 = 70$
    \item Se usan $3 \times 35{,}5 = 106{,}5$ L. Quedan $255{,}5 - 106{,}5 = 149$ L.
    \item $100 - 99{,}9 = 0{,}1$; $+ 88{,}8 = 88{,}9$; $- 77{,}7 = 11{,}2$; $+ 66{,}6 = 77{,}8$; $- 55{,}5 = 22{,}3$
  \end{enumerate}

  \subsection*{La Multiplicación (Ejercicios 45--64)}

  \begin{enumerate}[leftmargin=*, label=\textbf{\arabic*.}]
    \setcounter{enumi}{44}
    \item $92$
    \item $610$
    \item $132$
    \item $1{,}5$
    \item $2$
    \item $5$
    \item $1\,250 \times 345 = 431\,250$
    \item $4\,350 \times 128 = 556\,800$
    \item $6\,250 \times 164 = 1\,025\,000$
    \item $14\,250 \times 35 = 498\,750$
    \item $2\,250 \times 140 = 315\,000$
    \item $240 \times 13{,}25 = 3\,180$ Bs.
    \item $1\,420{,}5 \times 35{,}4 = 50\,285{,}7$
    \item $650{,}8 \times 42{,}5 = 27\,659$
    \item $4\,500 \times 37{,}5 = 168\,750$
    \item $12\,500 \times 348 = 4\,350\,000$
    \item $450{,}5 \times 230{,}4 = 103\,795{,}2$ m$^2$
    \item $0{,}125 \times 48\,000 = 6\,000$; \quad $0{,}0625 \times 32\,000 = 2\,000$
    \item $1\,200 \times 8{,}50 = 10\,200$ Bs.
    \item a) $(250 \times 40) \times 125 = 10\,000 \times 125 = 1\,250\,000$. \quad b) $125{,}50 \times 480 = 60\,240$ Bs.
  \end{enumerate}

  \subsection*{La Unidad Seguida de Ceros (Ejercicios 65--84)}

  \begin{enumerate}[leftmargin=*, label=\textbf{\arabic*.}]
    \setcounter{enumi}{64}
    \item $70$
    \item $340$
    \item $500$
    \item $25$
    \item $475$
    \item $8$
    \item $7$
    \item $0{,}35$
    \item $8$
    \item $125$
    \item $56\,000$
    \item $45$
    \item $1{,}25$
    \item $25$
    \item $0{,}07$
    \item $4{,}8$
    \item $0{,}85 \div 100 = 0{,}0085$ kg
    \item $4$
    \item $1\,500$
    \item a) $2{,}5 \times 1000 = 2\,500$; \quad $2\,500 \div 100 = 25$. \quad b) Hay que dividir $5$ veces: $32\,000 \to 3\,200 \to 320 \to 32 \to 3{,}2 \to 0{,}32$.
  \end{enumerate}

  \subsection*{La División (Ejercicios 85--104)}

  \begin{enumerate}[leftmargin=*, label=\textbf{\arabic*.}]
    \setcounter{enumi}{84}
    \item $21$
    \item $97 \div 3 = 32{,}33$ (resto $0{,}01$)
    \item $25$
    \item $19 \div 3 = 6{,}33$ (resto $0{,}01$)
    \item $7{,}3 \div 3 = 2{,}43$ (resto $0{,}01$)
    \item $1{,}2$
    \item $156\,250 \div 25 = 6\,250$
    \item $486\,500 \div 12 = 40\,541{,}66$ (resto $0{,}08$)
    \item $125{,}4 \div 0{,}7 = 1254 \div 7 = 179{,}14$ (resto $0{,}02$ en $1254 \div 7$)
    \item $1\,562\,500 \div 125 = 12\,500$
    \item $812{,}5 \div 0{,}3 = 8125 \div 3 = 2\,708{,}33$ (resto $0{,}01$ en $8125 \div 3$)
    \item $435\,000 \div 150 = 2\,900$
    \item $15\,625\,000 \div 1\,250 = 12\,500$
    \item $7\,680\,000 \div 3\,200 = 2\,400$
    \item $8\,125\,000 \div 2\,500 = 3\,250$
    \item $15\,625 \div 0{,}125 = 125\,000$
    \item $450 \div 0{,}7 = 4500 \div 7 = 642{,}85$ (resto $0{,}05$ en $4500 \div 7$)
    \item $27{,}60 \div 3 = 9{,}20$ Bs.
    \item $25{,}5 \div 7 = 3{,}64$ m por vestido (resto $0{,}02$ m)
    \item a) $(3\,600 \div 45) = 80$; $(0{,}25 \times 16) = 4$; $80 \div 4 = 20$. \quad b) Si $x \div 0{,}25 = 40$, entonces $x = 40 \times 0{,}25 = 10$.
  \end{enumerate}

  \vspace{10pt}
  \begin{cuadrosolucion}
    \textbf{\textquestiondown Cómo te fue?} Si llegaste hasta el final, ya eres todo un campeón de la aritmética. Recuerda: los errores no son fracasos, son pistas para aprender. \textquestiondown Hasta la guía 2!
  \end{cuadrosolucion}

\fi
\end{document}
